% =============================================================================
%  二、问题分析
% =============================================================================
\section{问题分析}

\subsection{三个映射：从问卷作答到状态轨道}

把一个“学习策略推荐”问题写成数学问题，本质上要连续做三次映射，每一次映射
都会丢失一部分信息，也都需要一次假设来补足。

\textbf{映射一（量化）：作答 $\to$ 投入产出向量。}设第 $i$ 名学生回答了
$p=14$ 个题项，第 $j$ 题的原始作答为 $r_{ij}$。不同题型的取值域不同（五级
有序、0—12 计数、0—20 小时连续），不能直接相加。记量化算子为
$q_j:\mathcal{R}_j\to[0,100]$，则个体被表示为
\begin{equation}
\bx_i=\bigl(x_{i1},x_{i2},x_{i3},x_{i4}\bigr)^{\top}\in\R_+^4,\qquad
y_i\in[0,100],
\label{eq:map1}
\end{equation}
其中 $x_{i1}$ 为听说读背等输入型训练的合成得分，$x_{i2}$ 为考试与模拟考的
训练量，$x_{i3}$ 为复习、辅导与课外拓展的时间投入，$x_{i4}$ 为心理与休息的
调节水平，$y_i$ 为期末成绩。量化算子的选择决定了后续一切结论的尺度，因此它
是复现工作的第一道关口。

\textbf{映射二（评价）：向量 $\to$ 效率。}在给定参照集 $\mathcal{S}$（同一
类学生的集合）下，用投入导向的 CCR 模型把个体映射为一个标量
\begin{equation}
\theta_i=\min\Bigl\{\theta\;\Big|\;\exists\,\bm{\lambda}\ge \bm{0}:\;
\sum_{j\in\mathcal{S}}\lambda_j\bx_j\le\theta\bx_i,\;
\sum_{j\in\mathcal{S}}\lambda_j y_j\ge y_i\Bigr\}\in(0,1].
\label{eq:map2}
\end{equation}
$\theta_i=1$ 表示位于前沿，$\theta_i<1$ 表示在同等产出下投入可压缩的比例。
效率把四维投入压缩成一维比值，代价是丢失了“哪一项投入最冗余”的信息——这
一信息只能从对偶乘子或松弛变量中部分找回。

\textbf{映射三（分层）：效率 $\to$ 层级。}把连续效率离散化为有限个有序等
级，是引入时间维度的前提。原文采用剔除式分层：先取 $\theta=1$ 的强有效单
元为第 1 层并从集合中剔除，在剩余集合上重算效率，再取新出现的 $\theta=1$
单元为第 2 层，如此反复。设第 $k$ 层集合为 $E_k$，则
\begin{equation}
E_1=\bigl\{i:\theta_i^{\mathcal{S}}=1\bigr\},\quad
E_{k}=\Bigl\{i:\theta_i^{\mathcal{S}\setminus\cup_{j<k}E_j}=1\Bigr\},\ k\ge 2 .
\label{eq:map3}
\end{equation}
这一映射有一个不易察觉的性质：层号不是个体的固有属性，而是“相对于当前参照
集的位置”。参照集一变，层号就变。这一点在第二节已作为局限提出，在第三节
将给出另一种口径（类内效率分位带），它把层号变成“同类群体中的百分位”，从
而保证了状态空间的占满性与可比性。

三次映射复合起来，一名学生在一个学期内就对应一条轨道
\begin{equation}
\bigl(r_{i1},\dots,r_{ip}\bigr)\;\longrightarrow\;\bx_i,y_i
\;\longrightarrow\;\theta_i\;\longrightarrow\;S_i\in\{1,\dots,5\}.
\label{eq:track}
\end{equation}
原文止步于 $S_i$，把它当作静态标签；本文则把 $S_i$ 当作随机过程的取值，让
它随学期内的时间演化。

\subsection{静态评价的数学结构与三处缺口}

把 DEA 写成式\eqref{eq:map2} 的形式后，可以清楚地看到它的三个数学特征，这
三个特征正是静态评价的边界。

\textbf{其一，评价是“截面”的。}$\theta_i$ 只依赖于同一时点的
$\{\bx_j,y_j\}_{j\in\mathcal{S}}$，与时间无关。若把 $\theta$ 视为一个随机
变量，DEA 给出的是它的一个实现值，而不是它的分布或演化规律。

\textbf{其二，评价是“相对”的。}$\theta_i$ 的大小取决于参照集的构成：当参
照集中出现一个投入更少、产出更高的个体时，所有人的效率都会下降。这意味着
效率值不可跨群体加总，也不可跨学期直接比较；原文的分层因此只能在样本内部
解释。

\textbf{其三，评价是“无记忆”的。}$\theta_i$ 不携带历史信息：一名学生上一
学期效率是 0.5、本学期是 0.8，与另一名一直稳定在 0.8 的学生，在静态评价中
完全等价。教育干预真正关心的“进步速度”“退步风险”“维持时间”在静态框架中
没有位置。

由此可以列出三处缺口，它们分别对应本文要用随机过程补上的三件事：

\begin{enumerate}[label=（\arabic*）,itemsep=2pt,topsep=2pt,leftmargin=3em]
\item \textbf{无转移概率}：给定当前层级 $k$，下一阶段处于各层级的概率未知。
  补法是把层级序列建成连续时间马尔可夫链，用转移强度矩阵 $\bQ$ 描述“上
  一层”与“下一层”的瞬时速率。
\item \textbf{无时间尺度}：从第 $k$ 层到第 1 层的平均等待时间未知。补法是
  求首达时间的期望 $m_k=\E_k[T_1]$，它满足一个三对角线性方程组，可解析求
  解并用蒙特卡洛校验。
\item \textbf{无策略空间}：原文的“路径”是同伴的复制，不是可优化的决策。补
  法是把行为改变写成动作集 $\mathcal{A}$，在学时预算约束下最大化
  $\E[\theta_T]$ 或最小化 $\E[T_1]$，从而得到“做哪几项、各做多少”的答案。
\end{enumerate}

\subsection{为什么随机过程是合适的语言}

学习效率的演化有三个特征，恰好与随机过程的三类基本对象对应。

\textbf{（1）状态是离散、有序、可反复穿越的。}层级只能取 1—5，且 1 优于 2
优于 $\cdots$ 优于 5；学生可以在相邻层之间来回移动，不能一次跨两层。这与
生灭过程（birth-death process）的状态结构完全一致：状态空间
$\{1,\dots,K\}$，只允许相邻转移，$1$ 与 $K$ 为反射壁。把“上行”记为出生、
“下行”记为死亡，则过程由两列强度刻画：
\begin{equation}
\lambda_k=\lim_{\Delta t\downarrow 0}\frac{\Pbb(S_{t+\Delta t}=k-1\mid S_t=k)}{\Delta t},
\qquad
\mu_k=\lim_{\Delta t\downarrow 0}\frac{\Pbb(S_{t+\Delta t}=k+1\mid S_t=k)}{\Delta t},
\label{eq:rates}
\end{equation}
生成矩阵 $\bQ=(q_{kl})$ 满足 $q_{k,k-1}=\lambda_k$、$q_{k,k+1}=\mu_k$、
$q_{kk}=-(\lambda_k+\mu_k)$。这一结构的最大好处是：平稳分布有闭式解，首达
时间有线性方程组，二者都可以写成可复算的公式，而不是黑箱输出。

\textbf{（2）强度依赖于协变量，个体之间存在系统性差异。}把 $\lambda_k,\mu_k$
写成行为协变量 $\bz_i$ 的函数，可以让“哪一种行为改变最能提高上行强度”成为
可估计的参数问题。这一层对应 Cox 型强度模型与多状态模型的极大似然估计。

\textbf{（3）效率本身是连续的，离散化会损失信息。}层级是粗粒化，真实效率
$\theta_t$ 在 $[0,1]$ 上连续演化，受大量未观测因素（教师风格、教材难度、
同伴影响）驱动。把它写成均值回复的扩散过程
\begin{equation}
\dd\theta_t=\kappa\bigl(\theta^{*}-\theta_t\bigr)\dd t+\sigma\dd W_t,
\qquad \kappa>0,
\label{eq:ou}
\end{equation}
既能刻画“效率有回归均值的天性”，又能给出首达某一目标水平的期望时间闭式
解，从而把“需要多久”的问题从离散链推广到连续时间。

三层结构（离散状态层、强度协变量层、连续扩散层）并非重复，而是分别回答
“在哪一层”“多快能上一层”“连续效率如何波动”。第三节将依次实现这三层。

\subsection{数据可识别性与建模假设}

本文最需要在方法论上说清楚的一点是：\textbf{没有纵向面板数据时，转移强度
不能被单期截面唯一识别。}设观测只有一期的层级 $S_i$，则 $\bQ$ 的
$K(K-1)$ 个参数中，只有平稳分布 $\bpi$ 可被截面频率约束，其余自由度需要
额外信息。本文采用三条互补的信息源来“部分识别”：

\begin{enumerate}[label=（\arabic*）,itemsep=2pt,topsep=2pt,leftmargin=3em]
\item \textbf{协变量差异}：不同学生在复习、应试、输入训练、心理作息与兴趣
  拓展上的差异，使 $\rho_i=\exp(\alpha+\bdelta^{\top}\bz_i)$ 可以被识别到
  相对大小；个体间 $\rho$ 的排序由数据决定，绝对尺度不可识别。
\item \textbf{时间锚点}：以原文给出的“S480 一个学期单步提升 25.38\%”为锚，
  反解时间单位，得到 $c$（每周的总体强度水平）。这是唯一的外部信息，也是
  本文唯一“借来”的参数。
\item \textbf{平稳性约束}：把模型隐含的平稳分布与观测到的层级频率对比，
  用拟合优度检验判断模型是否至少能复现截面结构。
\end{enumerate}

在此框架下，本文明确列出五条建模假设，后续所有结论都以此为前提：

\begin{description}[leftmargin=2em,itemsep=1pt,topsep=2pt]
\item[A1] 层级有序且相邻可达，1 为最优、5 为最差，两端为反射壁；
\item[A2] 在一个学期内，给定协变量后层级演化是时间齐次的马尔可夫链（行为
  习惯的惯性远大于其变化速度）；
\item[A3] 个体之间条件独立，即不存在同伴效应的即时传染（这一条在结论部分
  作为局限讨论）；
\item[A4] 层内效率同质，用层内均值 $\bar\theta_k$ 作为该层的代表效率，用于
  把层级轨道折算为效率轨道；
\item[A5] 行为改变的时间尺度以“周”为单位，一个学期按 16 周计。
\end{description}

\subsection{技术路线与记号约定}

据此，全文的计算路线为：先用原文口径复现 $\theta_i$ 与层级，得到状态序列；
再以状态序列与协变量估计 $\rho_i$ 与 $c$，得到 $\bQ_i$；用 $\bQ_i$ 计算
转移概率、期望效率与首达时间；用 OU 扩散补充连续效率的时间尺度；最后在
动作集上做预算约束优化。为避免符号混淆，本文统一约定：

\begin{itemize}[leftmargin=2em,itemsep=1pt,topsep=2pt]
\item $i=1,\dots,n$ 表示个体，$n=519$；$j$ 表示投入指标，$j=1,\dots,4$；
\item $k,l\in\{1,\dots,5\}$ 表示层级，$k=1$ 为最高效层；
\item $\theta_i$ 为类内 CCR 效率，$S_i$ 为层级，$\bar\theta_k$ 为层均效率；
\item $\Omega_i,P_i,D_i,H_i$ 为原文三准则（吸引力、进步度、邻近度、障碍度）；
\item $\lambda_k,\mu_k$ 为生灭过程的上行、下行强度，$\bQ$ 为生成矩阵，
  $\bP(t)=\exp(\bQ t)$ 为转移概率矩阵；
\item $T_1=\inf\{t:S_t=1\}$ 为首达最高层的时间，$m_k=\E_k[T_1]$。
\end{itemize}
